% Source: https://artofproblemsolving.com/community/c3h1504581_help_on_mathcounts_problem
% Stable identity: aops-topic-1504581-q01
Let $N=3^{2011}$ be the sum of the $k$ consecutive positive integers $a,a+1,\dots,a+k-1$, with $a\ge1$. Then
\[
N=ka+\dfrac{k(k-1)}{2},\qquad\text{so}\qquad 2N=k\,(2a+k-1).
\]
Write $m=2a+k-1$. Then $km=2\cdot3^{2011}$ and $m\ge k+1>k$. Also $m-k=2a-1$ is odd, so $k$ and $m$ have opposite parity. Conversely, any factorisation $2N=km$ with $m>k$ and $m-k$ odd gives a valid sum, with $a=\dfrac{m-k+1}{2}\ge1$.

The only factor $2$ goes to one of $k$, $m$, so $\{k,m\}=\{3^i,\,2\cdot3^j\}$ with $i+j=2011$. From $k<m$ and $km=2\cdot3^{2011}$ we get
\[
k^2<2\cdot3^{2011},\qquad k<3^{1005}\sqrt6\approx2.45\cdot3^{1005}.
\]
The possible values of $k$ are $3^i$ and $2\cdot3^j$. The largest of them below this bound is $k=2\cdot3^{1005}$, since $3^{1006}=3\cdot3^{1005}$ is already too large. It works: $m=3^{1006}$ is odd and greater than $k$, and
\[
a=\dfrac{3^{1006}-2\cdot3^{1005}+1}{2}=\dfrac{3^{1005}+1}{2}
\]
is a positive integer.

The last digits of the powers of $3$ repeat with period $4$: $3,9,7,1$. Since $1005\equiv1\pmod4$, $3^{1005}$ ends in $3$, so $k=2\cdot3^{1005}$ ends in $\boxed{6}$.
